\chapter{Adaptive authenticity}
\label{chap:adaptive}

\textbf{Status.} No originality is claimed for this chapter. We expect a bound of this kind to be
known to specialists, but we have not found a published statement of it (see below). Everything in it is
formalized: the reduction to the bad event (Section~\ref{sec:adaptive-lean}) and the bound on
the bad event (Section~\ref{sec:adaptive-bad}).

\textbf{Where the problem comes from.} In the protocol of \cite{JKO13} the problem does not
arise: the OT is the functionality $\mathcal{F}_{\mathrm{COT}}$, which fixes the prover's choice bits
before the garbled circuit is sent (\cite[Fig.~3 and Fig.~6]{JKO13}), and the authenticity needed
is static \cite[Definition~2]{JKO13}. It arises only when $\mathcal{F}_{\mathrm{COT}}$ is replaced by
an OT from which the receiver's choice cannot be extracted when its message is sent, as with the
Chou--Orlandi protocol of Chapter~\ref{chap:ot}, whose receiver message $B$ does not fix the
choice. As Byali, Patra, Ravi and Sarkar \cite[Sect.~3]{BPRS17} and Doerner, Kondi, Lee and shelat
\cite[App.~A]{DKLs18} point out, a corrupted receiver can delay the decryption, so a simulator
cannot extract its choice before the end of the protocol. Chou and Orlandi acknowledge the issue in
the revised version of their paper \cite[Sect.~1.1]{CO15}, where the attacking party is called the
sender, and note that it matters, for instance, when the OT is combined with garbled circuits. They
withdraw the claim of UC security and suggest, as a
remedy, a receiver message that proves timely decryption. This is a different issue from the one raised by Genç, Iovino and Rial \cite{GIR}, who
point out a flaw in the UC security proof of \cite{CO15} for a corrupted \emph{sender}; a claim of
theirs about a corrupted receiver, made in an earlier version, was withdrawn \cite[Sect.~1]{GIR}. An OT that realizes $\mathcal{F}_{\mathrm{COT}}$ avoids
it; Chapter~\ref{chap:pvw} does this with the OT of \cite{PVW08}.

With the Chou--Orlandi OT the receiver's message $B$ does not commit to a choice: the prover
receives both masked labels of every input together with the material, and it decides which key
point of $H_K$ to query, that is, which label to unmask, after seeing the material. Soundness with
the real OT therefore needs authenticity when the input $x$ is chosen \emph{adaptively}, after the
material and possibly after queries to $H$. The bound of Chapter~\ref{chap:random-oracle} is for a
fixed $x$, and its proof writes the material through the ideal garbler for that $x$
(Theorem~\ref{thm:real-ideal}), which is circular when $x$ is chosen later. Bellare, Hoang and
Rogaway \cite{BHR12} define adaptive authenticity in two forms \cite[Sect.~4, Fig.~6]{BHR12}:
$\mathrm{aut1}$, where the adversary chooses $x$ after seeing the garbled circuit, and
$\mathrm{aut2}$, where the bits of $x$ are chosen one at a time and each one is answered with its
token. The game below is $\mathrm{aut2}$, with access to the random oracle added. Their results
(Theorems~5 and~6, proved in their full version) are transforms that turn any statically secure
scheme into an adaptively secure one, by masking the garbled circuit until the input is given;
they do not say whether privacy-free Free-XOR garbling itself is $\mathrm{aut2}$-secure, which is
what we would need here. More recently, Barnum, Heath, Kolesnikov and Ostrovsky prove adaptive
security of unmodified schemes, half-gates among them, in the non-programmable random oracle model
\cite{BHKO26}; but their notion is adaptive \emph{privacy} \cite[Definition~5]{BHKO26} (and
obliviousness), not authenticity, and their argument (\cite[Theorem~9]{BHKO26}: query hiding plus
private rekeying) is of the same kind as the bad event of Section~\ref{sec:adaptive-bad}. Guo, Yang, Wang, Yu and Liu prove
that unmodified half-gates (the version with privacy, with a hash built from a random permutation) is
adaptively secure in the non-programmable random permutation model \cite[Sect.~5,
Theorem~1]{GYWYL}; their notion is again simulation-based privacy \cite[Definition~1]{GYWYL},
and their method, the H-coefficient technique, bounds the statistical distance between the real and
the simulated transcripts. We have not found a published
statement of adaptive authenticity for privacy-free garbling.

\section{The game}

\begin{definition}[Adaptive authenticity game]
  \label{def:adaptive-game}
  \lean{GcZk.Adaptive.realO, GcZk.Adaptive.realRun, GcZk.Adaptive.RealWin}
  \leanok
  \uses{def:garble, def:eval}
  The garbler samples $\Delta \in V \setminus \{0\}$, the zero-labels of the inputs and of the
  constants, and a random oracle $H$, and sends the material $\Mat$. The adversary then interleaves,
  in any order chosen adaptively:
  \begin{itemize}
    \item queries $(i, y)$ to $H$;
    \item label requests: for an input $j$ not yet requested and a bit $b$ of its choice, it
    receives the label $L_j + b\Delta$.
  \end{itemize}
  It takes $q$ actions in all, queries and requests together. In the formalization a repeated
  request of an input is answered with the label given the first time, whatever the bit. At the
  end the requested bits form $x$ (unrequested bits are $0$) and the adversary outputs a
  label $Z$. It wins if $Z = X_{\mathrm{out}} + \Delta$, where $X_{\mathrm{out}}$ is the active
  output label for $x$.
\end{definition}

If the adversary has requested every input that the circuit reads, it can compute
$X_{\mathrm{out}}$ from the material, the requested labels and the queries $(i, S_i)$ at the active
selector labels, and then winning amounts to finding $\Delta$. In general it must produce
$X_{\mathrm{out}} + \Delta$, and an unrequested input read by the circuit only makes this harder.

\begin{remark}[The model]
  \label{rem:adaptive-model}
  The hash is a uniformly random function $H$ and the adversary is computationally unbounded; only
  its number $q$ of actions is bounded. Authenticity is a game, not a simulation-based notion: no
  simulator answers the adversary's queries, so the oracle is never programmed. The simulated game
  of Definition~\ref{def:adaptive-sim} and the maps between randomness of
  Lemmas~\ref{lem:adaptive-to-sim} and \ref{lem:adaptive-fresh} are devices of the proof, which
  counts the values of $(\Delta, L, L_c, H)$ for which the adversary wins; the bound is a statement
  about the real game with $H$ uniform. The distinction between programmable and non-programmable
  random oracles, which matters for the simulation-based privacy results of \cite{BHKO26} and
  \cite{GYWYL}, does not arise here. What the bound does assume is the random oracle model itself,
  with $H$ used only by the garbling scheme; instantiations of $H$ are not covered (see
  Chapter~\ref{chap:limitations}), and in a protocol that uses a hash elsewhere the garbling must
  hash in a separate domain (Remark~\ref{rem:adaptive-plan}). Circuits are lists of nodes in evaluation order, and every
  operand refers to an earlier position; ``position'' and ``largest $i$'' below refer to this order.
\end{remark}

\section{The simulated game}

\begin{definition}[Simulated game]
  \label{def:adaptive-sim}
  \lean{GcZk.Adaptive.simM, GcZk.Adaptive.simO, GcZk.Adaptive.simRun, GcZk.Adaptive.SimWin}
  \leanok
  \uses{def:adaptive-game}
  The simulator sends uniform ciphertexts $T_i$ and uniform active constant labels, answers every
  label request with a fresh uniform label $R_j$ whatever the bit, and answers queries with a
  random function $H'$. It never uses $\Delta$; we sample $\Delta$ at the end. All active labels
  are defined by evaluation: for an input $x$, the active label of an input $j$ is $R_j$, and the
  evaluator reads $H'$ at the active selector labels.
\end{definition}

\section{The real game without bad queries}
\label{sec:adaptive-lean}

The zero-labels $L_s$ of the simulated game of Definition~\ref{def:adaptive-sim} depend on the final
choice $x$, so two runs with different completions do not share them. This breaks a direct extension
of the injection of Chapter~\ref{chap:random-oracle}: the map from real to simulated randomness has
to know $x$ to move the hash values at the inactive points, and on a run with a bad query $x$ cannot
be read back from the image. We therefore take the bad event in the \emph{real} game, where the two
points $L_s$, $L_s + \Delta$ of every selector are fixed by the garbler's randomness, and for the
final $x$.

\begin{definition}[The adaptive game in Lean]
  \label{def:adaptive-lean}
  \lean{GcZk.Adaptive.AAdv, GcZk.Adaptive.RealWin, GcZk.Adaptive.RealBad, GcZk.Adaptive.SimWin}
  \leanok
  \uses{def:adaptive-game, def:adaptive-sim}
  The adversary is deterministic (a randomized one is a family indexed by its coins) and takes $q$
  actions, each a query $(i, y)$ to $H$ or a request $(j, b)$ for an input $j < m$. Only the first
  request of an input counts: later requests of the same input are answered with the label of the
  first bit. The real randomness is $(L, L_c, H, J)$: the zero-labels of the inputs and of the
  constants, the hash and a junk component, which the real game does not use. A query $(i, y)$ is
  \emph{bad for the final input} if $y = S_i + \Delta$, where $S_i$ is the active label, for the
  final $x$, of the selector of the multiplexer at position $i$. The simulated game is the one of
  Definition~\ref{def:adaptive-sim}, with randomness $(R, A, H', T)$: requested labels, active
  constant labels, hash and ciphertexts.
\end{definition}

Section~\ref{sec:adaptive-bad} bounds the probability of a query bad for the final input.

\begin{lemma}[Identical until bad]
  \label{lem:adaptive-to-sim}
  \lean{GcZk.Adaptive.toAdSim, GcZk.Adaptive.simM_toAdSim, GcZk.Adaptive.simRun_toAdSim,
    GcZk.Adaptive.simWin_toAdSim}
  \leanok
  \uses{def:adaptive-lean, lem:to-sim}
  For $\Delta \neq 0$, map $(L, L_c, H, J)$ to $(R, A, H', T)$, where $R_j = L_j + x_j \Delta$ for
  the final $x$, $A$ are the active constant labels, $H'$ is $H$ with its values at the inactive
  points of the final $x$ replaced by $J$ (Lemma~\ref{lem:to-sim}), and $T_i$ is the real
  ciphertext of node $i$ if it is a multiplexer and the pad $J_i$ of Lemma~\ref{lem:to-sim}
  otherwise; the injectivity below uses both. The simulated material of the image is the real material. Without a bad query, the
  simulated game on the image runs exactly as the real game, and a real win is a simulated win.
\end{lemma}

\begin{proof}
  \leanok
  Requests are answered with $R_j$ in both games, since the first requested bit is the final one.
  Queries are answered by $H$ and $H'$, which differ only at the inactive points. The active output
  label is the same with $H$ and with $H'$, because the evaluator reads the hash only at active
  points.
\end{proof}

\begin{lemma}[Injectivity]
  \label{lem:adaptive-inj}
  \lean{GcZk.Adaptive.toAdSim_injOn, GcZk.Adaptive.idealAux_pad_inj}
  \leanok
  \uses{lem:adaptive-to-sim}
  For $\Delta \neq 0$ the map of Lemma~\ref{lem:adaptive-to-sim} is injective on the real
  randomness without a bad query.
\end{lemma}

\begin{proof}
  \leanok
  The image determines the run of the simulated game, which is the real run, hence the final $x$.
  Then $R$ and $A$ give $L$ and $L_c$. The material determines the pads of
  Lemma~\ref{lem:to-sim} at the multiplexers, and $T$ gives them elsewhere; with $H'$, injectivity
  of the map of Lemma~\ref{lem:to-sim} gives $(H, J)$.
\end{proof}

\begin{theorem}[Winning without a bad query]
  \label{thm:adaptive-not-bad}
  \lean{GcZk.Adaptive.win_not_bad_bound, GcZk.Adaptive.win_le}
  \leanok
  \uses{lem:adaptive-inj}
  Let $V$ be finite. With $\Delta$ uniform on $V \setminus \{0\}$ and the zero-labels and the hash
  uniform, the adversary wins without a bad query with probability at most $1/(|V| - 1)$. Hence it
  wins with probability at most $1/(|V| - 1)$ plus the probability of a bad query.
\end{theorem}

\begin{proof}
  \leanok
  For each $\Delta \neq 0$, Lemmas~\ref{lem:adaptive-to-sim} and \ref{lem:adaptive-inj} inject the
  real wins without a bad query into the simulated wins. The simulated game does not use $\Delta$,
  so each simulated randomness wins for at most one $\Delta$. Sum over $\Delta$ and swap the sums.
\end{proof}

With Theorem~\ref{thm:adaptive-not-bad}, Theorem~\ref{thm:adaptive} reduces to
Theorem~\ref{thm:adaptive-bad-bound}, the bound on bad queries.

\section{Bad queries}
\label{sec:adaptive-bad}

We bound the probability of a query that is bad for the final input. The difficulty is that this
event is not decided at the time of the query: whether a point is inactive depends on bits the
adversary requests later. We replace it by an event that is decided at the time of the query, and
for that we need to know which selector labels the adversary can know at that time. A selector
whose \emph{value} is fixed by the bits requested so far does not qualify: the XOR of two wires
that always carry opposite values has a fixed value, but its label may involve hash values the
adversary has never seen. The right notion is a combinatorial ``linear part'' of each label.

Fix the adversary and a step $k$. Let $R_k$ be the set of inputs requested before step $k$, and
$x^{(k)}$ the input formed by the bits requested before step $k$, with $0$ elsewhere. For a wire
$w$, $v_w(x)$ is its value on the input $x$ and $X^{(k)}_w$ its active label for $x^{(k)}$; in the
real game $X^{(k)}_w = Z_w + v_w(x^{(k)}) \Delta$, where $Z_w$ is the zero-label. For a
multiplexer at position $i$, $s_i$ is its selector and $c_i$ the candidate selected by
$v_{s_i}(x^{(k)})$.

\begin{definition}[Linear part]
  \label{def:lin}
  \lean{GcZk.Lin.pv, GcZk.Lin.pl, GcZk.Lin.plain_getD}
  \leanok
  \uses{def:adaptive-game}
  The \emph{atoms} are the symbols $a_j$, for the inputs $j$, and $b_i$, for the multiplexers $i$.
  The linear part $\ell^{(k)}_w \in \F^{\mathrm{atoms}}$ of a wire is defined by induction:
  \begin{itemize}
    \item an input $j$: $a_j$ if $j \notin R_k$, and $0$ if $j \in R_k$; a constant: $0$;
    \item an XOR $w = a \oplus b$: $\ell^{(k)}_a + \ell^{(k)}_b$;
    \item a multiplexer $i$: $\ell^{(k)}_{c_i}$ if $\ell^{(k)}_{s_i} = 0$, and
    $b_i + \ell^{(k)}_{c_i}$ otherwise.
  \end{itemize}
  A wire is \emph{known} at step $k$ if $\ell^{(k)}_w = 0$, and \emph{unknown} otherwise. The
  linear part depends only on the circuit, $R_k$ and $x^{(k)}$, not on the randomness.
\end{definition}

The atom $b_i$ stands for the hash value $H(i, X^{(k)}_{s_i})$ read by the evaluator at an unknown
selector, and $a_j$ for the label of an input not yet requested. Lemma~\ref{lem:adaptive-shift}
makes this precise.

\begin{lemma}[Known wires are determined]
  \label{lem:lin-det}
  \lean{GcZk.Lin.pv_add, GcZk.Lin.pv_eq_of_known}
  \leanok
  \uses{def:lin}
  Let $x'$ agree with $x^{(k)}$ on $R_k$. There is $\varphi : \mathrm{atoms} \to \F$ with
  $v_w(x') + v_w(x^{(k)}) = \sum_\beta \ell^{(k)}_w(\beta) \varphi(\beta)$ for every wire $w$. In
  particular a known wire has the same value, and in the real game the same active label, on $x'$
  and on $x^{(k)}$.
\end{lemma}

\begin{proof}
  \leanok
  Put $\varphi(a_j) = x'_j + x^{(k)}_j$ and $\varphi(b_i) = v_i(x') + v_i(x^{(k)}) + v_{c_i}(x') +
  v_{c_i}(x^{(k)})$, and argue by induction on the position. Inputs, constants and XORs are
  immediate. At a multiplexer with an unknown selector the claim is the definition of
  $\varphi(b_i)$ together with the induction hypothesis for $c_i$. At a multiplexer with a known
  selector the selector has the same value on $x'$ and on $x^{(k)}$ by induction, so both select
  $c_i$, and $v_i = v_{c_i}$ on both inputs. The active label is $Z_w + v_w \Delta$.
\end{proof}

\begin{lemma}[Monotonicity]
  \label{lem:lin-mono}
  \lean{GcZk.Lin.pl_mono, GcZk.Lin.known_mono}
  \leanok
  \uses{def:lin, lem:lin-det}
  For $j \le k$ there is a linear map $\psi$ with $\ell^{(k)}_w = \psi(\ell^{(j)}_w)$ for every
  wire $w$. In particular a wire known at step $j$ is known at step $k$.
\end{lemma}

\begin{proof}
  \leanok
  Put $\psi(a_t) = a_t$ if $t \notin R_k$ and $0$ otherwise, and $\psi(b_i) = \ell^{(k)}_i +
  \ell^{(k)}_{c^{(j)}_i}$, where $c^{(j)}_i$ is the candidate selected at step $j$. By induction on
  the position: XORs are linear; at a multiplexer with a selector unknown at step $j$,
  $\psi(\ell^{(j)}_i) = \psi(b_i) + \psi(\ell^{(j)}_{c^{(j)}_i}) = \ell^{(k)}_i +
  \ell^{(k)}_{c^{(j)}_i} + \ell^{(k)}_{c^{(j)}_i} = \ell^{(k)}_i$. At a multiplexer with a selector known at step $j$, the selector is known at step
  $k$ (its linear part is $\psi(0) = 0$), and by Lemma~\ref{lem:lin-det}, since $x^{(k)}$ agrees with
  $x^{(j)}$ on $R_j$, it has the same value at both steps; so the same candidate $c$ is selected,
  and $\ell^{(k)}_i = \ell^{(k)}_c = \psi(\ell^{(j)}_c) = \psi(\ell^{(j)}_i)$.
\end{proof}

\begin{definition}[Bad step]
  \label{def:adaptive-bad}
  \lean{GcZk.Adaptive.BadStep}
  \leanok
  \uses{def:lin}
  Step $k$ of the real game is \emph{bad} if its action is a query $(i, y)$ at a multiplexer $i$
  with selector $s$, and either $s$ is known at step $k$ and $y = X^{(k)}_s + \Delta$, or $s$ is
  unknown at step $k$ and $y \in \{Z_s, Z_s + \Delta\} = \{X^{(k)}_s, X^{(k)}_s + \Delta\}$.
\end{definition}

\begin{lemma}[First bad step]
  \label{lem:adaptive-first-bad}
  \lean{GcZk.Adaptive.exists_badStep, GcZk.Adaptive.not_realBad_of_first,
    GcZk.Adaptive.realMux_some}
  \leanok
  \uses{def:adaptive-bad, def:adaptive-lean, lem:lin-det}
  If some query is bad for the final input, some step is bad. If $k$ is the first bad step, the game
  truncated to its first $k$ steps has no query bad for its final input, which is $x^{(k)}$.
\end{lemma}

\begin{proof}
  \leanok
  Let the query at step $k$ be bad for the final input $x$: $y = X_s + \Delta$ for the active label
  $X_s$ of $s$ on $x$. Since $x$ agrees with $x^{(k)}$ on $R_k$, if $s$ is known at step $k$ then
  $X_s = X^{(k)}_s$ by Lemma~\ref{lem:lin-det}; otherwise $y$ is one of $Z_s, Z_s + \Delta$ anyway.
  For the second claim, let $j < k$ and let the query at step $j$ be $(i, y)$ with selector $s$.
  If $s$ is known at step $j$, then $X^{(k)}_s = X^{(j)}_s$ by Lemma~\ref{lem:lin-det}, and $y \ne
  X^{(j)}_s + \Delta$ because step $j$ is not bad. If $s$ is unknown at step $j$, $y$ is neither of
  its two points.
\end{proof}

\begin{lemma}[The image of a first bad step]
  \label{lem:adaptive-bad-image}
  \lean{GcZk.Adaptive.simBadAt_toAdSim, GcZk.Adaptive.simX_toAdSim, GcZk.Adaptive.SimBadAt}
  \leanok
  \uses{lem:adaptive-first-bad, lem:adaptive-to-sim, lem:adaptive-inj, lem:lin-mono}
  Let $\Delta \neq 0$ and $k < q$. The map of Lemma~\ref{lem:adaptive-to-sim} for the game truncated
  to $k$ steps is injective on the real randomness whose first bad step is $k$. On its image the
  simulated game, for the same $k$ steps and the same action at step $k$, has: the action at step
  $k$ is a query $(i, y)$ at a multiplexer, and, with $X$ the simulated active label of its selector
  $s_i$ for $x^{(k)}$, either $y = X + \Delta$, or $s_i$ is unknown at step $k$, $y = X$, and no
  query before step $k$ is at a point $(i', X_{s_{i'}})$ of a multiplexer $i'$ whose selector is
  unknown at step $k$.
\end{lemma}

\begin{proof}
  \leanok
  Injectivity is Lemma~\ref{lem:adaptive-inj} for the truncated game, which has no bad query by
  Lemma~\ref{lem:adaptive-first-bad}. The simulated run on the image coincides with the real run for
  $k$ steps (Lemma~\ref{lem:adaptive-to-sim}), so the action at step $k$ is the same, and the active
  labels for $x^{(k)}$ are the same in both games, because the evaluator reads the hash only at
  active points. A selector unknown at step $k$ is unknown at every earlier step
  (Lemma~\ref{lem:lin-mono}), so a query at one of its points before step $k$ would be a bad step.
\end{proof}

\begin{lemma}[Shifting an atom]
  \label{lem:adaptive-shift}
  \lean{GcZk.Adaptive.plab_shift, GcZk.Adaptive.plab_shiftH, GcZk.Adaptive.plab_shiftR,
    GcZk.Adaptive.simEval_getD, GcZk.Lin.pl_support}
  \leanok
  \uses{def:lin, def:adaptive-sim}
  In the simulated game at step $k$, give the atoms the values $\alpha(a_j) = R_j$ and $\alpha(b_i) =
  H'(i, X_{s_i})$, the labels being those for $x^{(k)}$. Let $s$ be an unknown wire, and let
  $\beta^*$ be $b_{i^*}$ for the largest $i^*$ with $b_{i^*}$ in $\ell^{(k)}_s$ if there is one,
  and some $a_j$ in $\ell^{(k)}_s$ otherwise. Add $v$ to the randomness at $\beta^*$: to $R_j$, or to
  $H'$ at the point $(i^*, X_{s_{i^*}})$. Then the labels of the wires before position $i^*$ do not
  change, and $X_s$ changes by exactly $v$.
\end{lemma}

\begin{proof}
  \leanok
  Write $\delta$ for the change. By induction on the position, $\delta X_w = \sum_\beta
  \ell^{(k)}_w(\beta)\, \delta\alpha(\beta)$. XORs are linear. At a multiplexer with an unknown
  selector, $X_i = \alpha(b_i) + X_{c_i}$, plus $T_i$ if the selector is $1$. At a multiplexer with a
  known selector, $X_i = H'(i, X_{s_i}) + X_{c_i}$ (plus $T_i$); the selector label does not change
  (its linear part is $0$), and neither does $H'$ at gate $i$, since $i \neq i^*$. The atoms change
  as follows: $\delta\alpha(\beta^*) = v$; the other $a_j$ do not change; $\delta\alpha(b_i) = 0$ for
  $i < i^*$, because the labels before $i^*$ do not change (by the same induction, the linear parts
  of those wires involve only such atoms). Since $\ell^{(k)}_s$ contains no $b_i$ with $i > i^*$,
  $\delta X_s = v$.
\end{proof}

\begin{lemma}[Fresh labels]
  \label{lem:adaptive-fresh}
  \lean{GcZk.Adaptive.card_freshHit_le, GcZk.Adaptive.shiftAt_injOn,
    GcZk.Adaptive.simRun_shiftAt}
  \leanok
  \uses{lem:adaptive-shift}
  For each $k$, the simulated randomness for which the action at step $k$ is a query $(i, y)$ with
  $s_i$ unknown at step $k$, $y = X_{s_i}$, and no earlier query at the active label of a selector
  unknown at step $k$ (the condition of Lemma~\ref{lem:adaptive-bad-image}; in particular none at
  the point of $\beta^*$ for $s = s_i$, when $\beta^*$ is some $b_{i^*}$), is at most a fraction
  $1/|V|$ of all of it.
\end{lemma}

\begin{proof}
  \leanok
  Map a pair $(\sigma, v)$, with $\sigma$ in this set and $v \in V$, to $\sigma$ shifted by $v$ at
  $\beta^*$. The shift does not change the view up to step $k$: $R_j$ is not requested before step
  $k$, and $H'$ changes at a point not queried before step $k$. So the image determines the run up
  to step $k$, hence $x^{(k)}$, the linear parts, $\beta^*$ and its point (the labels before $i^*$
  do not change). Then $v = X_{s_i} - y$ on the image, by Lemma~\ref{lem:adaptive-shift}, and
  $\sigma$ is recovered. So the map is injective, and the set has at most $|\mathrm{Sim}| / |V|$
  elements.
\end{proof}

\begin{theorem}[Bad queries]
  \label{thm:adaptive-bad-bound}
  \lean{GcZk.Adaptive.card_bad_le, GcZk.Adaptive.card_bad_le_first, GcZk.Adaptive.card_first_le,
    GcZk.Adaptive.sum_hitDelta_le}
  \leanok
  \uses{lem:adaptive-bad-image, lem:adaptive-fresh}
  With $\Delta$ uniform on $V \setminus \{0\}$ and the rest uniform, a query bad for the final input
  occurs with probability at most $q \bigl( 1/(|V| - 1) + 1/|V| \bigr)$, for an adversary taking
  $q$ actions.
\end{theorem}

\begin{proof}
  \leanok
  By Lemma~\ref{lem:adaptive-first-bad} it suffices to bound, for each $k < q$, the randomness whose
  first bad step is $k$. Fix $\Delta \neq 0$ and inject it into the simulated randomness by
  Lemma~\ref{lem:adaptive-bad-image}. The images with $y = X + \Delta$ occur, for each simulated
  randomness, for at most one $\Delta$, since the simulated game does not use $\Delta$. The others
  lie, for every $\Delta$, in the set of Lemma~\ref{lem:adaptive-fresh}. Summing over $\Delta$, the
  count is at most $|\mathrm{Sim}| + (|V| - 1)|\mathrm{Sim}|/|V|$; divide by $(|V| - 1)
  |\mathrm{Sim}|$ and sum over $k$.
\end{proof}

\begin{theorem}[Adaptive authenticity]
  \label{thm:adaptive}
  \lean{GcZk.Adaptive.adaptive_authenticity, GcZk.Adaptive.card_win_le}
  \leanok
  \uses{thm:adaptive-not-bad, thm:adaptive-bad-bound}
  In the adaptive authenticity game, an adversary taking $q$ actions wins with probability at most
  \[ \frac{q + 1}{|V| - 1} + \frac{q}{|V|} . \]
\end{theorem}

For $|V| \geq 2$ this is at most $(2q + 1)/(|V| - 1)$, an elementary inequality that is not part of
the formal statement.

\begin{proof}
  \leanok
  Theorem~\ref{thm:adaptive-not-bad} and Theorem~\ref{thm:adaptive-bad-bound}.
\end{proof}

\begin{remark}[The Chou--Orlandi OT]
  \label{rem:adaptive-plan}
  Theorem~\ref{thm:adaptive} is used in Chapter~\ref{chap:co-soundness}, with
  Lemma~\ref{lem:keys-sub} (both labels of an input require solving square Diffie--Hellman) and a
  simulation of the OT layer by label requests, to extend the soundness of
  Chapter~\ref{chap:commitment} to the Chou--Orlandi OT. The OT derives its keys with a hash as
  well. Theorem~\ref{thm:adaptive} is for an oracle used only by the garbling, so the two must hash
  in separate domains (for instance with distinct prefixes), which makes them independent random
  oracles.
\end{remark}
